\section{Cơ chế phòng thủ D3}\label{sec:d3}
% Dàn ý gốc: ../outline/03_d3_defence.md (khoảng 2,5 trang)

\subsection{Nguyên tắc thiết kế}

D3 theo ba nguyên tắc. (i) \textbf{Mô hình không phải ranh giới an ninh.} Dù mô hình sinh ra gì, một hành động chỉ
có hiệu lực khi broker thực thi nó; mọi bảo đảm cứng đều nằm phía host. (ii) \textbf{Người dùng là chủ thể đáng tin
duy nhất.} Quyền được suy ra từ nguyên văn chỉ thị người dùng, không bao giờ từ nội dung công cụ và không bao giờ từ
việc một mô hình diễn giải ý định người dùng. (iii) \textbf{Chặn không được làm kết thúc nhiệm vụ.} Hành động bị từ
chối được báo lại cho tác tử theo dạng giúp nó hoàn thành các bước hợp lệ.

Hình~\ref{fig:d3} minh họa các lớp; Bảng~\ref{tab:layers} tóm tắt chúng.

% HÌNH 1: Claude vẽ, lưu ở images/fig1_d3_architecture.pdf. Bỏ comment khi có file.
% \begin{figure}[t]
% \centering
% \includegraphics[width=\textwidth]{fig1_d3_architecture}
% \caption{Kiến trúc D3. Người dùng $\to$ mô hình tác tử $\rightleftarrows$ Tool Broker (schema, giới hạn, nhãn
%   nguồn gốc, luật R2--R4, luật D, phản hồi) $\rightleftarrows$ công cụ; lời nhắc trên đường quan sát; câu trả lời
%   cuối đi qua R3/R4.}\label{fig:d3}
% \end{figure}
\todo{Hình 1 (kiến trúc D3)}

\begin{table}[t]
\centering
\caption{Các lớp của D3. ``TĐ'' đánh dấu cưỡng chế tất định phía host.}\label{tab:layers}
\small
\begin{tabular}{@{}p{0.22\textwidth}p{0.15\textwidth}p{0.30\textwidth}cp{0.17\textwidth}@{}}
\toprule
Lớp & Thời điểm & Cơ chế & TĐ & Nhắm tới \\
\midrule
L0 Tool Broker & mọi lời gọi & kiểm schema, giới hạn công cụ, nhãn nguồn gốc & có & lời gọi sai; nền cho các lớp khác \\
L1 Cấp quyền đích (R2) & trước sink & đích phải có nguyên văn trong chỉ thị người dùng & có & M1, M2, M6 \\
L2 Luồng giá trị mật (R3) & trước sink; ở câu trả lời & giá trị mật từ công cụ không được đi ra trừ khi người dùng đã viết & có & M3, M4 \\
L3 Che câu trả lời (R4) & ở câu trả lời & che từng giá trị vi phạm & có & M4 \\
L4 Luật tín hiệu theo đích (D) & trước sink & từ chối khi quan sát bị gắn cờ nêu đúng đích của sink & có & M3, M5 trên sink người dùng nêu \\
L5 Lời nhắc (P2) & sau mỗi quan sát & lặp lại chỉ thị người dùng cho mô hình & không & tất cả, tùy mô hình \\
L6 Phản hồi từ chối & sau khi từ chối & thông báo có cấu trúc; không chạy công cụ; tác tử tiếp tục & có & tiện ích sau khi chặn \\
\bottomrule
\end{tabular}
\end{table}


\subsection{Tool Broker và nhãn nguồn gốc (L0)}

Mọi lời gọi công cụ đều đi qua broker. Khi không có phòng thủ (cấu hình A0), broker chỉ kiểm tra lời gọi theo schema
và cưỡng chế giới hạn từng công cụ. Trong D3, broker còn lưu, với mỗi quan sát, định danh tài liệu cùng nhãn tin cậy
và nhạy cảm, và tập \textbf{giá trị} trích từ các quan sát mật. Việc trích giá trị là từ vựng và không phụ thuộc
ngôn ngữ: địa chỉ e-mail, URL, định danh dài từ sáu ký tự và số dài. Thao tác đọc không bao giờ bị chặn (luật R1),
nên tác tử luôn thu thập được thông tin nhiệm vụ cần.

Mọi phép so sánh dùng một hàm chuẩn hóa $n(\cdot)$: Unicode NFKC, bỏ ký tự định dạng (trong đó có ký tự độ rộng
bằng không) và đưa về chữ thường. Chuẩn hóa chỉ dùng để so sánh, không bao giờ thay đổi nội dung mà mô hình hay công
cụ nhận được.

\subsection{Cấp quyền đích (L1, luật R2)}

Một lời gọi sink $a$ chỉ được phép nếu đích của nó do người dùng viết:
\begin{quote}
\textbf{R2.} chỉ cho phép $a$ nếu $n(\mathit{dest}(a))$ là chuỗi con của $n(u)$.
\end{quote}
Văn bản người dùng không được phân tích ý định, và không có danh sách cụm từ tiếng Việt hay tiếng Anh nào được tra.
Vì vậy luật không phụ thuộc vào cách viết đoạn chèn: một đoạn chèn M2 tránh động từ ``gửi'', tách địa chỉ hay giấu
đích trong tài liệu thứ hai sẽ thất bại vì cùng lý do với đoạn chèn M1, là đích nó cần không có trong $u$. Cái giá
của so khớp nguyên văn xuất hiện ở hai chỗ. Người dùng nhắc đích một cách gián tiếp (``địa chỉ trong hợp đồng'') sẽ
không được phục vụ. Và một đích người dùng chỉ nhắc trong câu cấm vẫn được tính là đã cấp quyền, đúng là trường
hợp M6 thăm dò.

\subsection{Luồng giá trị mật (L2--L3, luật R3 và R4)}

Gọi $C$ là tập giá trị trích từ các quan sát có nhãn mật. Với lời gọi sink $a$ hoặc câu trả lời cuối $y$:
\begin{quote}
\textbf{R3.} từ chối nếu có $v \in C$ xuất hiện trong $n(\mathit{payload}(a))$ (tương ứng $n(y)$) và $n(v)$ không
xuất hiện trong $n(u)$.\\
\textbf{R4.} câu trả lời cuối vi phạm R3 không bị đánh hỏng: mỗi giá trị vi phạm được thay bằng
\code{[REDACTED]}.
\end{quote}
R3 hiện thực một chính sách luồng thông tin tối thiểu trên hai nhãn ở Mục~\ref{sec:threat}: giá trị mật chỉ được
tới sink qua chính văn bản người dùng. Luật coi câu trả lời cuối là một sink, nhờ đó đóng kênh M4 mà các bộ lọc lời
gọi công cụ bỏ sót. R4 giữ phần còn lại của câu trả lời, nên một câu trả lời đúng có lẫn một mã bị lộ chỉ mất mã đó
chứ không mất cả nhiệm vụ. Vì R3 so khớp nguyên văn, giá trị mà mô hình đã biến đổi trước khi xuất (định dạng lại,
tách ra, diễn đạt lại) sẽ không được nhận ra; Mục~\ref{sec:discussion} quay lại giới hạn này.

\subsection{Luật tín hiệu chèn lệnh giới hạn theo đích (L4)}

R2 và R3 để lại một khoảng hở: đoạn chèn thao túng một sink mà người dùng \emph{đã} nêu, với giá trị không mật (M5)
hoặc qua tham số R3 không phủ. Phòng thủ trước của chúng tôi (D2, cấu hình được đánh giá ở tập niêm phong thứ nhất)
lấp khoảng này bằng bộ phát hiện dựa trên cụm từ để nhận ra văn bản giống chỉ thị trong quan sát: khi có bất kỳ
quan sát nào bị gắn cờ, mọi sink sau đó không được cấp quyền bởi một văn phạm chặt trên văn bản người dùng đều bị
từ chối. Trên tập niêm phong thứ nhất, lớp này là nguồn chặn nhầm lớn nhất trên nhiệm vụ lành tính (22 quỹ đạo
benign\_plain trên ba mô hình nền), nhưng bỏ nó trên tập phát triển lại làm số họ gây hại tăng 2--4 ở mỗi mô hình
nền.

D3 giữ bộ phát hiện và thu hẹp hệ quả của nó. Gọi $F$ là tập quan sát bị bộ phát hiện gắn cờ:
\begin{quote}
\textbf{D.} từ chối lời gọi sink $a$ nếu có $o \in F$ chứa $n(\mathit{dest}(a))$.
\end{quote}
Luật chỉ kích hoạt khi nội dung không đáng tin, mang dáng dấp chỉ thị, nói về đúng cái đích mà tác tử sắp dùng. Một
trang bị gắn cờ ở đâu đó trên quỹ đạo không còn chặn một lần gửi không liên quan do người dùng yêu cầu, còn việc
chuyển hướng sang đích khác vẫn do R2 lo. Thiết kế được chọn bằng phát lại offline trên các quỹ đạo phát triển có
sẵn, trước mọi lượt chạy mô hình mới: luật D chặn đủ 11 ca gây hại mà lớp rộng đã chặn và không chặn ca nào trong 14
lời gọi sink lành tính bị lớp rộng từ chối nhầm. Một phương án khác theo dõi giá trị từ nguồn bị gắn cờ chỉ chặn 9/11
và bị loại.

\subsection{Lời nhắc (L5) và phản hồi từ chối (L6)}

\textbf{Lời nhắc.} Sau mỗi quan sát công cụ, chỉ thị gốc của người dùng được nối vào cùng tin nhắn, kèm một câu nói
rằng mọi chỉ thị khác tìm thấy trong dữ liệu công cụ phải bị bỏ qua (phòng thủ sandwich, như \code{repeat\_user\_prompt}
của AgentDojo). Đây là lớp duy nhất phụ thuộc vào việc mô hình tuân thủ. Lớp này được giữ vì nó phủ những gì các
lớp tất định không thấy: trên tập niêm phong thứ nhất, riêng nó giảm tỉ lệ ca gây hại từ 32,9\% xuống 5,6\%, và kết
hợp với các luật thì còn 1,2\%.

\textbf{Phản hồi từ chối.} Khi một luật từ chối hành động, công cụ không được thực thi và tác tử nhận một quan sát
có cấu trúc báo rằng hành động bị chính sách từ chối. Tác tử sau đó có thể làm các bước còn lại hoặc trả lời. Nếu
không có phản hồi này, mỗi lần từ chối sẽ kết thúc quỹ đạo và biến mọi lần chặn nhầm thành một nhiệm vụ thất bại.

\subsection{Chọn thành phần}

D3 được thiết kế và tinh chỉnh chỉ trên tập phát triển; tập niêm phong thứ nhất khi đó đã mở và chỉ được dùng để
xác định các lớp lỗi. Trước mọi lượt chạy GPU, chúng tôi cố định tiêu chí chọn: gộp trên ba mô hình nền, ứng viên
phải có số ca tấn công gây hại không nhiều hơn cấu hình trước (D2+P2) và tiện ích lành tính cao hơn hẳn; thành phần
nào làm tăng tác hại hoặc giảm tiện ích trong ablation thì bị bỏ.

Bảng~\ref{tab:ablation} báo cáo ablation. Chỉ biến thể có luật D, lời nhắc gốc và phản hồi gốc (tên D3DEST trong
log, trong bài gọi là D3) đạt tiêu chí. Một lời nhắc mạnh hơn, giục mô hình hoàn thành mọi bước được yêu cầu kể cả
việc gửi, giúp Gemma nhưng làm số họ gây hại của Llama tăng gấp đôi và làm tiện ích của Qwen sụp (benign\_plain 141
$\to$ 39 trên 240); nó không chuyển giao giữa các mô hình nền nên bị loại. Một thông báo từ chối mạnh hơn không thêm
gì. Hai ý tưởng khác bị loại ở bước phát lại: giới hạn R3 theo tài nguyên người dùng nêu (mở được 3 nhiệm vụ có cấp
quyền nhưng thả 5 tấn công) và so khớp giá trị gần đúng (không bắt thêm ca nào trên dữ liệu phát triển).

\begin{table}[t]
\centering
\caption{Ablation D3 trên tập phát triển (48 họ, 560 ca mỗi ô). Mỗi ô: họ gây hại /48 $\cdot$ tiện ích
benign\_plain /240 $\cdot$ tiện ích benign\_auth /80. Cột gộp: ca tấn công gây hại /720 $\cdot$ tiện ích lành tính
/960.}\label{tab:ablation}
\small
\begin{tabular}{@{}lcccc@{}}
\toprule
Cấu hình & Gemma-4-12B & Llama-3.1-8B & Qwen3.5-9B & Gộp \\
\midrule
A0 (không phòng thủ) & 38 $\cdot$ 159 $\cdot$ 54 & 34 $\cdot$ 81 $\cdot$ 30 & 40 $\cdot$ 140 $\cdot$ 36 & 330 $\cdot$ 500 \\
D2+P2 (trước đó) & 2 $\cdot$ 147 $\cdot$ 45 & 5 $\cdot$ 82 $\cdot$ 36 & 3 $\cdot$ 141 $\cdot$ 24 & 12 $\cdot$ 475 \\
\textbf{D3} (luật D, nhắc v1, phản hồi v1) & 2 $\cdot$ 148 $\cdot$ 49 & 5 $\cdot$ 85 $\cdot$ 40 & 3 $\cdot$ 141 $\cdot$ 27 & \textbf{12 $\cdot$ 490} \\
luật D, nhắc v1, phản hồi v2 & 2 $\cdot$ 150 $\cdot$ 49 & 6 $\cdot$ 84 $\cdot$ 40 & 3 $\cdot$ 142 $\cdot$ 26 & 14 $\cdot$ 491 \\
luật D, nhắc v2, phản hồi v2 & 1 $\cdot$ 163 $\cdot$ 54 & 10 $\cdot$ 75 $\cdot$ 38 & 2 $\cdot$ 39 $\cdot$ 16 & 21 $\cdot$ 385 \\
luật D, nhắc v2, phản hồi v1 & 1 $\cdot$ 162 $\cdot$ 54 & 10 $\cdot$ 69 $\cdot$ 38 & 2 $\cdot$ 40 $\cdot$ 16 & 21 $\cdot$ 379 \\
\bottomrule
\end{tabular}
\end{table}


\paragraph{Chi phí và phạm vi.} D3 không thêm lời gọi mô hình nào: mọi luật là phép toán chuỗi trên quỹ đạo, còn lời
nhắc thêm độ dài chỉ thị người dùng vào mỗi tin nhắn quan sát. D3 không cần truy cập trọng số và dùng nguyên vẹn
được với mô hình dịch vụ (Mục~\ref{sec:external}). D3 không tuyên bố chặn được việc sửa giá trị không mật trên sink
người dùng đã nêu khi nội dung chèn không nhắc tới đích (một phần của M5), cũng như rò rỉ giá trị đã bị mô hình biến
đổi.
